% !TeX root = ../../hodge-classes.tex
\section{Hodge algebraicity on mixed powers}\label{ext:sec:main}
\subsection{The mixed-powers and individual-powers statements}
Put $M=\Ql(\zeta_{42})=\Ql(\zeta_{21})$. Embeddings of $M$ are labeled
by units modulo $42$. In this section, $B$ is the underlying complex
isogeny class of $A_B$ with a relabeled $M$-action, and $\extD$
is an additional abstract sixfold. The notation $\extD$ is unrelated
to the quadrilateral generator $D_t$ or its vertex, which is equivalent
to $B_t$. We use $\Psi$, rather than the geometric types $\Phi$ of Paper I,
for the following rank-one $M$ CM types:
\[
 \Psi_B=\{1,5,11,13,17,19\},\qquad
 \Psi_{\extD}=\{1,5,11,13,19,25\}.
\]
In terms of the geometric type $\Phi_B$ just specified, $\Psi_B=19\Phi_B$
modulo $42$. This relabeling changes the field marking used in the
character calculation, not the underlying complex isogeny class.
Here rank one means that rational first cohomology is a one-dimensional
$M$-vector space. All varieties and maps in this section are over
$\C$. Let $Y$ be the Jacobian of the smooth projective model of
$y^2=x^7-1$.

\begin{theorem}[Mixed powers]\label{ext:thm:mixed}
For all nonnegative integers $r,s,t$, every rational Hodge class on
$B^r\times {\extD}^s\times Y^t$ is algebraic, in every codimension.
The statement is invariant under replacing any factor by a complex
isogenous abelian variety.
\end{theorem}

\begin{corollary}[The packet fibers and the additional abstract type]
Let $A_A,A_B,A_C$ be the three certified packet fibers of Theorem~\externalref{packet}{thm:A}.
Then every rational Hodge class on
$A_A^a\times A_B^b\times A_C^c\times \extD^d$ is algebraic for all nonnegative
$a,b,c,d$.
\end{corollary}
\begin{proof}
The product is isogenous to $Y^{2a}\times B^{b+c}\times {\extD}^d$.
An isogeny and its rational inverse act by algebraic correspondences;
pullback and pushforward divided by its degree transport algebraicity.
The isogeny identifications are Proposition~\externalref{packet}{prop:knownCM}.
They concern the underlying varieties, not their prescribed polarizations
or RM markings. Equal rank-one CM types identify rational Hodge
structures, and commensurable full lattices give a complex isogeny after
clearing denominators; see also \cite[\S3.11, Prop.~3.12]{MilneCM2020}.
\end{proof}

\begin{theorem}[Individual powers for the fixed field]\label{ext:thm:individual}
If $X$ is a complex abelian sixfold with a rational $M$-action and
$\dim_M H^1(X,\Ql)=1$, then every rational Hodge class on every positive
power $X^n$ is algebraic.
\end{theorem}

Theorem~\ref{ext:thm:individual} ranges over eight CM-type orbits, while
Theorem~\ref{ext:thm:mixed} concerns the stated three isogeny classes.
The later Theorem~\ref{eight:thm:all} establishes algebraicity on arbitrary
mixed products of all eight rank-one $M$ types. The symbol $\extD$ denotes an abstract CM type here; no new curve
point or fourth quotient class is asserted. None of these statements
identifies a standard trace polarization with a principal polarization,
or proves literal membership in a prescribed principally polarized
moduli locus.

\subsection{The algebraic input classes}\label{ext:sec:inputs}
For a smooth projective $T$, write
\[
 \extHdg^p(T)=H^{2p}(T,\Ql)\cap H^{p,p}(T),\qquad
 \extAlg^p(T)=\operatorname{im}\bigl(\mathrm{CH}^p(T)_{\Ql}
                   \longrightarrow H^{2p}(T,\Ql)\bigr).
\]
For a splitting field $E$, algebraicity of a split eigenline means
containment in $\extAlg^p(T)\otimes_{\Ql}E$. It does not assert that an
individual vector of that line is rational.

\subsubsection{The $B/Y$ input}
Let $J_m=J(y^2=x^m-1)$ for odd $m$. Shioda's original
Example~6.3 proves the rational Hodge conjecture for $J_{21}$ in all
codimensions \cite[pp.~76--77]{Shioda1981}. It is a theorem about this
Jacobian, not an assertion about all its powers. Its factorization is
\[
 J_{21}\sim X_{21}\times J_7\times J_3,
\]
where $X_{21}$ is the primitive six-dimensional factor. The primitive
holomorphic character set modulo $21$ is $\{1,2,4,5,8,10\}$;
multiplication by $19$ gives the restriction of $\Psi_B$.
Thus $X_{21}\sim B$ by CM classification. Projection pullback and
restriction to the zero point of $J_3$, together with isogeny transport,
prove algebraicity of all Hodge classes on $B\times Y$. This also
supplies algebraicity on $B$ itself. The source audit is retained in
\path{experiments/mixed_products_v100/literature_scope_review.json}.

\subsubsection{An explicit split Weil polarization for ${\extD}$}
Write $\zeta=e^{2\pi i/42}$, $u=\zeta+\zeta^{-1}$,
$L=M^+$, and $K_7=\Ql(\sqrt{-7})$. Set
\[
 s_3=1+2\zeta^{14}=i\sqrt3,\qquad
 s_7=\sum_{a=1}^{6}\chi_{-7}(a)\zeta^{6a}=i\sqrt7,
 \qquad u_a=\zeta^a+\zeta^{-a},
\]
where the signs are fixed at the displayed complex embedding. Put
\begin{align*}
 b&=\sqrt{21}\,u_1u_{13}u_{17}
   =4u^5+3u^4-15u^3-13u^2+4u+3,\\
 \xi&=s_7b,\qquad
 f(u)=u^6+u^5-6u^4-6u^3+8u^2+8u+1=0.
\end{align*}
The equality $\sqrt{21}=-s_3s_7$ fixes its positive sign. The exact
sign data at the real embeddings $a=1,5,11,13,17,19$ are
\[
\begin{array}{c|rrrrrr}
 a&1&5&11&13&17&19\\\hline
 \epsilon_a\text{ for }\Psi_{\extD}&+&+&+&+&-&+\\
 \chi_{-7}(a)&+&-&+&-&-&-\\
 \operatorname{sign}\sigma_a(b)&+&-&+&-&+&-
\end{array}
\]
Thus $\operatorname{Im}\phi(\xi)>0$ for $\phi\in\Psi_{\extD}$.
The alternating form on $M$,
\[
 E_\xi(x,y)=\Tr_{M/\Ql}(\xi x\bar y),
\]
satisfies $E_\xi(x,Jx)>0$ for $x\ne0$ and the Riemann bilinear
relations. It is integral on $\mathcal O_M$; for any full lattice
realizing this rational CM structure a positive integer multiple clears
denominators. It therefore supplies a polarization, not necessarily a
principal one. Its $K_7$ signature is $(3,3)$ and
$E_\xi(kx,ky)=\Nm_{K_7/\Ql}(k)E_\xi(x,y)$.

Let $T(x,y)=\Tr_{M/K_7}(bx\bar y)$. On the $K_7$-basis
$1,u,\ldots,u^5$ its determinant is
\[
 \det T=\operatorname{disc}(f)\Nm_{L/\Ql}(b)
       =(3^3 7^5)(-21^3)=-3^6 7^8=-64827^2.
\]
Indeed each $u_a$ used above has norm $1$ and
$\Nm_{L/\Ql}(\sqrt{21})=-21^3$. The associated Hermitian form is
\[
 E_\xi(s_7x,y)+s_7E_\xi(x,y)=-14T(x,y).
\]
The sixth-power scalar is a $K_7$-norm, and $64827^2$ is the norm of
the rational element $64827$. Hence the Hermitian discriminant class
is $-1$ in $\Ql^*/\Nm_{K_7/\Ql}(K_7^*)$; clearing lattice denominators
preserves this class. The identities and signs have distinct exact
implementations in
\path{experiments/fixed_field_cm_census_v98/k7_polarization_primary.py}
and \path{independent/k7_split_polarization_v99.py}.

Markman's Theorem~1.5.1 applies to a polarized abelian sixfold of Weil
type for an imaginary quadratic field, with Hermitian discriminant
$-1$ \cite[\S1.1 and Thm.~1.5.1]{Markman2025}. The constructed compatible
polarization and balanced signature verify its hypotheses for ${\extD}$.
The theorem does not require a principal polarization. Consequently the
two-dimensional rational Weil space of ${\extD}$ is algebraic. In particular
it supplies both split lines of the relation
\[
 (1,-1,1,-1,-1,-1)
\]
in the six positive directions $1,5,11,13,17,19$ of ${\extD}$.

\subsubsection{Five $B/{\extD}$ quartets from one source partition}
For degree $42$ Fermat factors use triples $(10,12,20)$ for $B$ and
$(1,14,27)$ for ${\extD}$. The least-residue rule for $H^{1,0}$ in
\cite[eq.~(2.1), p.~68]{Shioda1981} gives exactly $\Psi_B,\Psi_{\extD}$.
The $B$ triple has effective conductor $21$ and the ${\extD}$ triple conductor
$42$; both have twelve character labels and the indicated six-dimensional
admissible factors. Equality of marked CM types transports them to $B,{\extD}$
by complex isogenies.

The invariant support
\[
 S=\{(B,1),(B,19),({\extD},31),({\extD},41)\}
\]
has its twelve Fermat residues partitioned into
\[
 (2,10,31,41),\quad(12,15,18,39),\quad(14,28),\quad(20,22).
\]
For every unit modulo $42$, each quadruple has least-residue sum $84$;
the pairs are opposite. These are blocks in Shioda's
$\mathcal B_{42}^{2}$ and $\mathcal B_{42}^{0}$. The supportwise
proof of Theorem~4.3, using Lemmas~4.1--4.2, makes the selected exterior
character line algebraic \cite[pp.~70--72]{Shioda1981}. No odd-degree
hypothesis is imposed in these statements. This is not an application
of the restricted hyperelliptic Theorem~5.4 at even degree.

The five $B/{\extD}$ quartets in the lattice basis below are precisely the
translates $25S,\allowbreak 29S,\allowbreak 5S,\allowbreak 31S,\allowbreak S$. Their negatives are translates by $-1$.
Thus all required quartet lines are algebraic. The complete codimension-two
check, beyond these five needed inputs, is
\path{independent/fermat_codim2_v100.py}; it independently verifies all
$78=66+12$ invariant four-line supports and their source partitions.

\subsection{From an integral relation basis to algebraic tensors}
The following lemma isolates the unbounded argument. It concerns a
particular collection of CM factors, not a general assertion that
pairwise Hodge algebraicity settles arbitrary products.

\begin{lemma}[Algebraic contractions]\label{ext:lem:contraction}
Let $F_1,\ldots,F_m$ be complex CM abelian varieties. Suppose their
first cohomologies split into distinct nonzero weight lines within
each factor, each with a unique opposite, over a number field $E$.
Retain the factor labels and one representative of each opposite pair,
and let $\mathcal K$ be the full integer lattice of relations among
their synchronized Hodge characters. Suppose $\mathcal K$ has an
integral generating set whose signed relation tensors, for both signs,
are algebraic after extension to $E$. Suppose also that the embedding
line projectors are induced by algebraic correspondences with coefficients
in $E$. Then every rational Hodge class on every mixed power of the
$F_i$ is algebraic.
\end{lemma}
\begin{proof}
First a known exterior generator can be moved to singleton $H^1$ slots.
For a factor $F$ of dimension $d$, the polynomial
\[
 \pi_1=\prod_{\substack{0\le k\le2d\\k\ne1}}
 \frac{[n]^*-n^k\mathrm{Id}}{n-n^k},\qquad n\ge2,
\]
is the rational $H^1$ projector, because $[n]^*$ acts as $n^k$ in
degree $k$. It is an algebraic correspondence. Pull a wedge of distinct
embedding lines back along addition $F^q\to F$. Apply $\pi_1$ on
each factor and the assigned embedding projector in each slot. Exactly
one assignment survives, up to sign, yielding its pure Kunneth tensor
on separate copies. For a generator on several original copies, perform
this operation copy by copy. All classes remain in $\extAlg\otimes E$.

Fix an ample rational divisor $\theta$ on $F$. On $Z\times F\times F$,
diagonal pullback, cup product with $\theta^{d-1}$, and pushforward
along $Z\times F\to Z$ act on singleton entries by
\[
 \eta\otimes x\otimes y\longmapsto
 \eta\int_F x\wedge y\wedge\theta^{d-1}.
\]
This is algebraic and lowers codimension by one. Hard Lefschetz and
Poincare duality make the displayed alternating pairing nondegenerate.
It is Hodge-group invariant. Distinct weights with unique opposites
therefore force a nonzero pairing on each opposite pair. The scalar
lies in $E$ and can be divided out. Only slots of the same marked factor
are contracted. No inverse Lefschetz operator is used.

For $z\in\mathcal K$, write an exact integer combination of the known
generators; replace a negative coefficient by copies of the known
negative generator. Tensorize on fresh slots and take external products.
If a direction occurs $P$ times positively and $N$ times negatively,
contract $\min(P,N)$ opposite pairs. The surviving multiplicity is
exactly the signed coordinate $P-N$ of $z$, and all contractions are
nonzero. Rescaling gives the desired pure tensor in $\extAlg\otimes E$.
This permits nonconformal cancellation by correspondences, without
squaring an exterior line or extracting a root of a cycle class.

Now take an invariant exterior monomial on a target mixed power. Set
aside opposite pairs, algebraic by Lefschetz $(1,1)$; the remaining
multiplicities give $z\in\mathcal K$. Construct its singleton tensor,
reinsert the pairs by external products, and pull back along the product
of diagonals assigning each slot to its intended target copy. The target
labels are distinct, so this yields the required nonzero wedge up to
sign. Temporary source powers need no separately assumed Hodge
conjecture: the needed classes have been constructed on them.

The zero-weight exterior monomials span the rational Hodge space after
splitting the CM torus. They all lie in $\extAlg\otimes E$. The opposite
inclusion holds because algebraic classes are Hodge; faithful scalar
extension proves equality of the rational spaces, degree by degree.
If a target factor is absent, its temporary slots have signed total zero
and are all contracted away. A product with every exponent zero is a point.
\end{proof}

\paragraph{The required embedding projectors.}
For the present factors take a primitive element $t$ in a full-rank
order of actual CM endomorphisms. On $H^1$ the polynomial
\[
 p_\sigma=\prod_{\tau\ne\sigma}
       \frac{t^*-\tau(t)\mathrm{Id}}{\sigma(t)-\tau(t)}
\]
projects onto the $\sigma$ line. Compose it with $\pi_1$. The
denominators are nonzero, and the result is an algebraic correspondence
with coefficients in $E$; no model of the variety over $E$ is asserted.
For $B,{\extD}$ there are twelve embeddings of $M$; for $Y$ there are six
embeddings of $K=\Ql(\zeta_7)$. Thus $\pi_1$ uses degrees $0$--$12$
for $B,{\extD}$ and $0$--$6$ for $Y$, and contraction uses respectively
$\theta^5$ and $\theta^2$.

\subsection{The synchronized nine-generator lattice}\label{ext:sec:lattice}
Give $Y$ the $K$ action $\zeta_7\mapsto\rho^4$, where
$\rho$ is induced by $x\mapsto\zeta_7x$. Its holomorphic type is
$4\{1,2,3\}=\{1,4,5\}$ modulo $7$. This is the relative marking
used in the reviewed $B/Y$ source. The sixfold census type
$\{1,5,11,19,25,29\}$ is its inverse image modulo $42$ and describes
$Y^2$; it must not be counted as twelve distinct $H^1$ lines on one $Y$.

Use direction labels $(1,5,11,13,17,19)$ on $B,{\extD}$ and $(1,5,11)$
on $Y$, restricting the latter modulo $7$ to $(1,5,4)$. For the
\emph{same} $g\in G=(\Z/42\Z)^*$ put
\[
 w_{F,a}(g)=
 \begin{cases}
  +1,&ga\in\Psi_F,\\-1,&ga\notin\Psi_F,
 \end{cases}
\]
using modulus $42$ for $F=B,{\extD}$ and modulus $7$, with
$\Psi_Y=\{1,4,5\}$, for $Y$. Let $W$ be this $12\times15$ matrix.
No independent Galois translation is applied to a single factor. Each
factor's signed weights are distinct and nonzero, with a unique opposite.

The Hodge torus of this product is the smallest rational subtorus
containing its Hodge circle. Equivalently, its cocharacter space is
generated by the simultaneous Galois translates of the Hodge-circle
cocharacter. Pairing an embedding character with these translates gives
the entries $w_{F,a}(g)$, up to a common harmless sign convention.
Consequently a signed multiplicity vector defines a trivial Hodge
character exactly when it lies in $\ker_{\Z}W$. This uses the same $g$
in every factor: separate translations would describe a different
relative CM marking. Repeated copies repeat the same characters. In
even cohomological degree, Hodge-circle invariance is precisely Hodge
type $(p,p)$, and splitting the torus over $M$ therefore identifies its
invariant exterior monomials with the scalar extension of rational
Hodge classes.

The following rows are the full nine relation generators, in the column
order $B(6)\mid {\extD}(6)\mid Y(3)$:
\[
\setlength{\arraycolsep}{3.5pt}
\begin{array}{rrrrrr|rrrrrr|rrr}
0&0&0&0&0&0&1&-1&1&-1&-1&-1&0&0&0\\
0&0&0&1&-1&0&0&0&0&0&1&1&0&0&0\\
0&1&0&-1&0&0&0&0&0&1&1&0&0&0&0\\
0&1&1&0&0&0&0&-1&0&-1&0&0&0&0&0\\
1&0&-1&0&0&0&0&-1&1&0&0&0&0&0&0\\
1&0&0&0&0&1&-1&0&-1&0&0&0&0&0&0\\
1&-1&-1&0&0&0&0&0&0&0&0&0&0&0&1\\
1&0&-1&0&0&1&0&0&0&0&0&0&-1&0&0\\
1&0&0&-1&0&0&0&0&0&0&0&0&0&-1&1
\end{array}
\tag{$\mathcal R$}\label{ext:eq:basis}
\]
\begin{lemma}[The full integer kernel]\label{ext:lem:lattice}
The rows in \eqref{ext:eq:basis} form a $\Z$-basis of $\ker_{\Z}W$.
\end{lemma}
\begin{proof}
Order $G$ increasingly. The last six rows of $W$ are the negatives of
the first six in reverse order. The first six rows, in one-based columns
$1,2,3,4,5,7$, have determinant $-32$, so $\operatorname{rank}W=6$.
Every displayed vector is in its kernel. Project them to the one-based
coordinates
\[
 6,8,9,10,11,12,13,14,15.
\]
The resulting $9\times9$ integer matrix has determinant $+1$.
These exact assertions, including an explicit integer inverse and full
coordinate multiplication, are certified in
\path{experiments/mixed_triple_v102/primary.py} and
\path{experiments/mixed_triple_v102/primary.json}, with a separately
reconstructed implementation and review listed in
Section~\ref{ext:sec:evidence}.

The nine vectors form a rational basis of the rank-nine kernel, and
the indicated projection restricts to an isomorphism onto $\Ql^9$.
For an integral kernel vector, unimodularity gives integer coefficients
from its projected coordinates. Subtracting this combination leaves a
kernel vector with zero projection, hence zero. This proves equality
with the full integer kernel, not just equal rational rank.
\end{proof}

\begin{proof}[Proof of Theorem~\ref{ext:thm:mixed}]
The first row of \eqref{ext:eq:basis} is the pure-${\extD}$ Weil relation supplied
by the split polarization and Markman's theorem. Rows two through six
are the five $B/{\extD}$ quartet translates in Section~\ref{ext:sec:inputs}.
The last three are invariant $B/Y$ quartets, algebraic by the one-product
Shioda input. Both signs of all nine rows are algebraic. The exact weight
checks give distinct nonzero weights and unique opposites within each
factor, and the projectors above meet the remaining hypothesis of
Lemma~\ref{ext:lem:contraction}. Apply Lemma~\ref{ext:lem:lattice} and then
Lemma~\ref{ext:lem:contraction}.

In particular no separate theorem for ${\extD}\times Y$ is assumed. Its
classes can be constructed using temporary $B$ slots which contract
away. Pairwise Hodge algebraicity alone would not prove this theorem;
the full integral index-one statement is the extra input.
\end{proof}

\subsection{The eight fixed-field types and individual powers}
The exact CM-type census chooses one embedding from each of six conjugate
pairs, so its complete domain has $2^6=64$ elements. There are eight
Galois orbits, with the following representatives. Rank denotes the
connected Hodge-group rank, not the Mumford--Tate rank.
\begin{center}
\small
\begin{tabular}{@{}lrrl@{}}
\toprule
Representative modulo42 & Orbit & Rank & Weight structure\\
\midrule
$1,5,11,13,17,19$ &12&5&one unit relation; $B$\\
$1,5,11,13,17,23$ &6&3&three independent pairs, multiplicity2\\
$1,5,11,13,19,25$ &12&5&one unit relation; ${\extD}$\\
$1,5,11,17,23,29$ &12&6&six independent pairs\\
$1,5,11,19,25,29$ &6&3&three independent pairs, multiplicity2\\
$1,5,11,23,25,29$ &12&6&six independent pairs\\
$1,11,23,25,29,37$ &2&1&one pair, multiplicity6\\
$1,13,19,25,31,37$ &2&1&one pair, multiplicity6\\
\bottomrule
\end{tabular}
\end{center}
The separate exact census implementations are
\path{experiments/fixed_field_cm_census_v98/primary.py} and
\path{experiments/fixed_field_cm_census_v98/independent.py}.
Their retained outputs and hostile reviews bind the finite coverage,
orbit data and character-relation calculations.

\begin{proof}[Proof of Theorem~\ref{ext:thm:individual}]
For each of the six independent-weight orbits, an invariant exterior
monomial on any power chooses equally many lines of each opposite weight.
Pair them to write it as a product of degree-two invariant monomials.
This argument includes repeated weight spaces and divisors mixing the
copies. Lefschetz $(1,1)$ and rational scalar descent prove algebraicity.

For the two rank-five orbits the relation lattice of the six directions
is generated respectively by
\[
 (1,-1,-1,1,-1,1),\qquad(1,-1,1,-1,-1,-1).
\]
Every coordinate is $\pm1$; rank five and a coordinate equal to one
show that these are the full integral relation lattices. The two signs
give the two split Weil lines on the original sixfold. They are algebraic
on $B$ by the $J_{21}$ factor input, and on ${\extD}$ by the explicit
$K_7$ polarization and Markman's theorem. Apply the one-factor case of
Lemma~\ref{ext:lem:contraction}, or equivalently distribute these six-line
wedges among copies using addition and the embedding projectors.

Finally Galois relabeling of the $M$ action and commensurability of the
rank-one lattices transport the statement to every variety of a type
in the indicated orbit. Positive integer scaling clears the integral
polarization denominators without changing the relevant norm class.
This proves the stated fixed-field, individual-powers scope.
\end{proof}

\subsection{Certificate boundaries and retained extensions}\label{ext:sec:evidence}
The argument just given uses one pure Weil input and eight quartets.
The shorter mixed-powers proof does not require the complete Graver
basis or the subsequent cloud source searches. Those remain separate,
reviewed results rather than implicit hypotheses.

\begin{center}
\small
\begin{tabular}{@{}p{.10\linewidth}p{.82\linewidth}@{}}
\toprule
Claim & Retained mathematical scope\\
\midrule
C45 & Individual powers of all eight rank-one $M$ sixfold type orbits.\\
C46 & All mixed powers of the three packet fibers, via $B$ and $J_7$.\\
C47 & Codimension two on all $B^r\times {\extD}^s$.\\
C48 & All codimensions on one $B\times {\extD}$.\\
C49 & Complete signed Graver-basis and finite generator reduction;
the reduction alone asserts no new algebraicity.\\
C50 & All mixed $B/{\extD}$ powers by an integral basis and contractions.\\
C51 & All mixed $B/{\extD}/J_7$ powers, with the packet corollary above.\\
C52 & All mixed products of the six degree-$21$ Fermat-type classes.\\
C53 & All mixed products of all eight rank-one $M$ sixfold classes.\\
C54 & The stated full-degree CM-subfield factors and all their mixed products.\\
\bottomrule
\end{tabular}
\end{center}

The V101 Graver computation gives $122$ vectors up to sign, with degrees
$4,6,8,10$ and counts $6,44,24,48$. Its completeness uses the separate
signed-graph proof of a coordinate bound two and degree bound twelve.
Of the $244$ signed patterns, $168$ repeat an embedding direction.
The bounded source search produced independently verified Fermat
witnesses for $144$ of those signed patterns; the other $24$ were not
settled by that search. The contraction theorem establishes their
algebraicity without changing those recorded search outcomes.

The following repository-relative paths locate the source-bound proof
records. They are provenance references, not external assertions of
specialist approval:
\begin{itemize}
\item \path{experiments/fixed_field_cm_census_v98/final_extension_exposition_review.json}
and \path{experiments/fixed_field_cm_census_v98/k7_polarization_hostile_review.json};
\item \path{experiments/mixed_products_v100/second_primitive_review.json},
\path{experiments/mixed_products_v100/graph_review.json}, and
\path{experiments/mixed_products_v100/second_primitive_codim3_review.json};
\item \path{experiments/mixed_powers_v101/graver_review.json},
\path{experiments/mixed_powers_v101/contraction_review.json}, and
\path{experiments/mixed_powers_v101/source_results_review.json};
\item \path{experiments/mixed_triple_v102/primary.py},
\path{experiments/mixed_triple_v102/primary.json},
\path{independent/mixed_triple_v102.py},
\path{experiments/mixed_triple_v102/independent.json}, and
\path{experiments/mixed_triple_v102/final_review.json}.
\end{itemize}
The independent V102 source and exact artifact hashes are bound by its
comparison and final checkpoint reviews. Finite arithmetic comparisons
support the stated integer identities; they do not replace the
contraction, polarization, source-theorem or rational-descent arguments.
No full fresh packet-certificate execution or consolidated release is
claimed by these extension records.

\section{Mixed products from six fixed-field CM-type orbits}
\label{six:sec:main}
Put $M=\Ql(\zeta_{42})=\Ql(\zeta_{21})$ and
$G=(\Z/42\Z)^\times$. As in Section~\ref{ext:sec:main}, types refer to
rational first cohomology: an embedding label $g$ belongs to the type
when its eigenline lies in $H^{1,0}$. We use the canonical types $\Psi_i$
in Table~\ref{six:tab:types}, retaining their indices in the eight-orbit
census. These are abstract field markings, not additional geometric
CM points on the McMullen curve.

\begin{theorem}[Six-orbit mixed products]\label{six:thm:products}
Let $X_1,\ldots,X_N$ be complex abelian sixfolds, each with an embedding
$M\hookrightarrow\End^0(X_j)$ and $\dim_M H^1(X_j,\Ql)=1$. Suppose each cohomological
CM type is a Galois translate of one of the six types in
Table~\ref{six:tab:types}. Then every rational Hodge class on
$X_1\times\cdots\times X_N$ is algebraic, in every codimension.
The factors may repeat, and their lattices may be chosen independently.
The empty product is allowed.
\end{theorem}

\begin{table}[htbp]
\centering\small
\setlength{\tabcolsep}{4pt}
\begin{tabular}{r l c r r r r}
\toprule
$i$ & $\Psi_i$ & $t\pmod{21}$ & $q$ & $\dim A_t$ & $e$ & $h$\\
\midrule
0 & $\{1,5,11,13,17,19\}$ & $(1,1,19)$ & 21 & 6 & 1 & 19\\
1 & $\{1,5,11,13,17,23\}$ & $(1,7,13)$ & 21 & 6 & 1 & 11\\
4 & $\{1,5,11,19,25,29\}$ & $(3,3,15)$ & 7 & 3 & 2 & 11\\
5 & $\{1,5,11,23,25,29\}$ & $(1,6,14)$ & 21 & 6 & 1 & 1\\
6 & $\{1,11,23,25,29,37\}$ & $(1,4,16)$ & 21 & 6 & 1 & 1\\
7 & $\{1,13,19,25,31,37\}$ & $(7,7,7)$ & 3 & 1 & 6 & 1\\
\bottomrule
\end{tabular}
\caption{Degree-$21$ Fermat presentations. Here $q$ is the character
order, $e=[M:\Ql(\zeta_q)]$, and $h$ transports the induced type to
$\Psi_i$. An admissible factor $A_t$ need not be simple.}
\label{six:tab:types}
\end{table}

\begin{proof}
Let $F_m$ be the smooth plane Fermat curve. Its Jacobian has admissible
rational character factors indexed by Galois orbits of triples
$t=(t_1,t_2,t_3)$ with nonzero entries and sum zero modulo $m$.
If $t$ has order $q$, its factor $A_t$ has dimension $\varphi(q)/2$
and rank-one rational first cohomology over $E=\Ql(\zeta_q)$;
no simplicity assumption is required
\cite[\S2, eqs.~(2.2)--(2.4) and Definition~2.1]{Shioda1981}.
We specify the holomorphic convention used for these factors.
On $x^m+y^m=1$, the differentials
\[
 x^{a-1}y^{b-m}\,dx,
 \qquad a,b\geq1,\quad a+b\leq m-1,
\]
are regular and have diagonal character $(a,b,m-a-b)$.
At $y=0$, use $dx=-y^{m-1}x^{1-m}dy$; at infinity their local
factor has exponent $m-a-b-1\geq0$. Their distinct characters and
number $(m-1)(m-2)/2$ give a holomorphic basis. Thus the holomorphic
labels of $t$ are exactly those $u$ for which
$\sum_j\langle ut_j\rangle_m=m$, where residues lie in
$\{1,\ldots,m-1\}$, in agreement with
\cite[eq.~(2.1), p.~68]{Shioda1981}.

For each row of Table~\ref{six:tab:types}, set
$d=\gcd(21,t_1,t_2,t_3)$, $q=21/d$, and $v=t/d$.
The ordered character $v$ has image $\mu_q$. Its factor therefore
has the $E$-action just described, with holomorphic type
\[
 T_q(v)=\left\{u\in(\Z/q\Z)^\times:
              \sum_j\langle uv_j\rangle_q=q\right\}.
\]
The induced $M$-type is
\[
 T=\{g\in G:g\bmod q\in T_q(v)\}
   =\left\{g\in G:\sum_j\langle gt_j\rangle_{21}=21\right\}.
\]
Direct residue calculations give $hT=\Psi_i$ in each displayed row.
For example the embedding fibers have size $e=12/\varphi(q)$,
so the dimensions and copy multiplicities in the table follow
without identifying an elliptic or three-dimensional factor with
a sixfold. The type-stabilizer orders in row order are
$1,2,2,1,6,6$; in particular several types are induced.

We explain why induction gives an actual field action on $A_t^e$.
Write $V=H^1(A_t,\Ql)$. Choosing an $E$-basis of $M$ identifies
\[
 H^1(A_t^e,\Ql)=V^e\simeq M\otimes_E V.
\]
Multiplication on the first factor is a faithful rank-one $M$-action.
It preserves the Hodge decomposition: after splitting $E$, a matrix
over $E$ only mixes copies of the same embedding component, whose
Hodge type is homogeneous. The line labeled $g$ lies over the
$E$-embedding labeled $g\bmod q$, so its type is exactly $T$.

These rational Hodge endomorphisms are rational algebraic endomorphisms.
Indeed, dualizing a rational Hodge map to homology gives a real linear
map commuting with the complex structures. Clearing denominators
makes it a map of lattices, hence a holomorphic homomorphism of
complex tori. Its graph is algebraic for abelian varieties.
The reversal of composition under pullback causes no problem for
the commutative field $M$. Thus the construction supplies an embedding
$M\hookrightarrow\End^0(A_t^e)$, not just a choice of split eigenvalues.
Re-mark this action by $\iota_{\rm new}=\iota_{\rm old}\circ\sigma_h$.
An old label $g$ becomes $gh$, since
$\sigma_g(\sigma_h(a))=\sigma_{gh}(a)$; hence the new type is $hT$.

Two rank-one rational $M$ Hodge structures with the same type are
$M$-linearly Hodge-isomorphic: an $M$-linear isomorphism preserves each
embedding line and its Hodge bidegree. The weight-one lattice argument
above turns such an isomorphism into a complex algebraic isogeny after
clearing denominators. This works for induced types as well as simple
ones. A further field re-marking realizes any Galois translate.
Consequently each $X_j$ is complex-isogenous to the corresponding
$A_t^e$. Every finite product in the theorem is therefore a degree-$21$
Fermat-type abelian variety in Shioda's sense, up to isogeny.

The remaining input is the full semigroup condition $P_{21}$.
The sentence preceding the published Theorem~3.3 of da Silva,
with its Appendix computation, establishes $P_m$ for $m=21,27$
\cite{daSilva2021Fermat}. We use this published result as a literature
input, without claiming an independent replay of its computation.
The superscript in the accompanying hypersurfaces $X_m^n$ denotes
dimension, not Cartesian power. Shioda's Theorem~4.4 states that
$P_m^{6c-2}$ implies rational Hodge algebraicity in codimension $c$
for every degree-$m$ Fermat-type abelian variety
\cite[Theorem~4.4, p.~72]{Shioda1981}.
Full $P_{21}$ supplies this hypothesis for every $c\geq1$.
It applies to the identified product, including repeated admissible
factors; codimension zero is immediate. This use of Theorem~4.4
supplies the correspondence and multiplicity statement directly,
without inferring mixed algebraicity from separate factors.
Finally, isogeny pullback and pushforward preserve cycle classes,
with $f_*f^*=(\deg f)\mathrm{id}$, and transfer the conclusion to
the original product over $\Ql$.
\end{proof}

The finite identifications, including character kernels, embedding fibers,
and action transporters, are retained in
\path{experiments/six_orbit_v109/identification.json}; the human proof
and its primary-source hypotheses remain essential.
Theorem~\ref{six:thm:products} concerns the six listed isogeny classes.
The two remaining census types, $2$ and $3$, are addressed in
Section~\ref{eight:sec:main}. The six-class theorem imposes no principal polarization and
makes no claim of a new literal McMullen point or prescribed moduli
normalization. The equality $\Ql(\zeta_{42})=\Ql(\zeta_{21})$ does not
by itself lower the order of a degree-$42$ Fermat character.

\section{Arbitrary mixed products for the fixed degree-twelve field}
\label{eight:sec:main}
The six-class result and two additional Fermat-character cycles extend
algebraicity to arbitrary mixed products for the fixed field
$M=\Ql(\zeta_{42})$. This field restriction is essential.
\begin{theorem}[All mixed products for the fixed field]\label{eight:thm:all}
Let $X_1,\ldots,X_N$ be complex abelian sixfolds equipped with embeddings
$M\hookrightarrow\End^0(X_j)$ such that $\dim_M H^1(X_j,\Ql)=1$ for every $j$.
Then every rational Hodge class on $\prod_jX_j$ is algebraic,
in every codimension. Repetitions, arbitrary choices of lattices and
the empty product are allowed.
\end{theorem}

Choose representatives $X_i$ for the eight canonical CM-type orbits.
Besides the six types in Table~\ref{six:tab:types}, these are
\[
 \Psi_2=\{1,5,11,13,19,25\}=\Psi_{\extD},\qquad
 \Psi_3=\{1,5,11,17,23,29\}.
\]
The eight orbits exhaust the $2^6=64$ choices of one embedding from
each conjugate pair; their respective sizes are
$12,6,12,12,6,12,2,2$. The marked weight-one isogeny argument in the
proof of Theorem~\ref{six:thm:products} reduces the theorem to arbitrary
powers of these eight representatives, including induced types.

\subsection{Actual embedding lines and an integral relation certificate}
\label{eight:sec:lattice}
Set $R=(1,5,11,13,17,19)$ and order $G=(\Z/42\Z)^\times$ increasingly.
For every factor use the \emph{same} row $g\in G$ and set
\[
 w_{i,a}(g)=2\mathbf1_{\Psi_i}(ga)-1.
\]
The Hodge-torus argument of Section~\ref{ext:sec:lattice} shows that
zero sums of these columns are exactly the invariant signed tensor
characters. All twelve actual $M$-embedding lines on each $X_i$ remain
distinct even if their Hodge-torus weights coincide.
There are $48$ actual conjugate-pair coordinates $(i,a)$, $a\in R$.

Retain the following representatives for distinct pairs of opposite
weight columns, with the factors ordered by $i$ and the labels by size:
\[
\begin{array}{c|cccccccc}
i&0&1&2&3&4&5&6&7\\\hline
R_i&R&\{1,5,11\}&R&R&\{1,5,11\}&R&\{1\}&\{1\}.
\end{array}
\]
Let $W$ be the resulting $12\times32$ matrix and $e_{i,a}$ its
coordinate vectors. Every original column is uniquely a signed retained
column of the same factor. Thus there is a surjection
$P:\Z^{48}\to\Z^{32}$ and a section $J$ including the retained actual
labels, with $PJ=1$ and $W_{\rm act}=WP$. Consequently
\[
 \ker_{\Z}W_{\rm act}
 =J(\ker_{\Z}W)\oplus\ker P.
\]
The sixteen generators of $\ker P$ have the form
$f_{i,a}-\epsilon f_{i,b}$, where $a$ is not retained and
$w_{i,a}=\epsilon w_{i,b}$. Their tensors use two singleton copies:
$L_{i,a}\otimes L_{i,-b}$ if $\epsilon=1$, and
$L_{i,a}\otimes L_{i,b}$ if $\epsilon=-1$.
They are split components of rational degree-two Hodge classes and are
algebraic by Lefschetz $(1,1)$. This justifies the reduction without
identifying different geometric eigenlines or assuming they are
interchanged by a prescribed endomorphism.

Here is a precise finite certificate for the remaining integer kernel.
All coordinate indices in this paragraph are \emph{zero-based}.
The pivot columns of $W$ are $(0,1,2,3,4,6)$ and its rank is six.
Write $F$ for the increasing list of the other $26$ indices.
The unique kernel chart $K$ whose rows in $F$ form the identity
has integral entries. Thus $\ker_{\Z}W=K\Z^{26}$.
For $S=\{0,1,4,5,6,7\}$, use the same column order on $W_S$;
its analogous free-coordinate kernel chart has fourteen columns.
Embed these fourteen columns in $\Z^{32}$, and follow them by the
nine vectors in \eqref{ext:eq:basis}, embedding $B,{\extD},Y$
respectively in factors $0,2,4$ with the displayed retained labels.
Call this ordered list of $23$ known vectors $\mathcal B$.
They have rank $19$ and are algebraic by
Theorems~\ref{six:thm:products} and~\ref{ext:thm:mixed}; here
$X_4$ is isogenous to $Y^2$.

Define two further relations
\begin{align*}
 q&=-e_{0,1}-e_{0,17}+e_{1,1}+e_{2,1},\\
 r&= e_{0,1}-e_{0,5}+e_{0,13}-2e_{1,1}+e_{3,1}.
\end{align*}
For $h\in G$, send each label $a$ to $ha$ and then reduce to its
signed representative using $P$. Order the distinct resulting integer
vectors lexicographically in the fixed $32$ coordinates. Each of
$q$ and $r$ has twelve such translates. Form a square matrix by taking,
in order, the projections to $F$ of:
\begin{enumerate}[label=(\roman*)]
\item the first nineteen vectors of $\mathcal B$ (the fourteen
$S$-chart columns followed by the first five rows of
\eqref{ext:eq:basis});
\item the lexicographically first translate of $q$;
\item the first six lexicographically ordered translates of $r$.
\end{enumerate}
Its determinant is $-1$. Its retained integer inverse, multiplied by
the corresponding full $32$-coordinate matrix, gives exactly $K$.
These explicitly specified finite identities are certified in
\path{experiments/all_eight_v109/lattice.json}, reconstructed by
\path{experiments/all_eight_v109/review/independent_check.py}, and
regressed in \path{tests/test_v109_reductions.py}.
The chart identity proves integral surjectivity onto the free
coordinates, while unimodularity proves that the known vectors and
the two seed orbits generate the full integer kernel. Equal rational
rank alone is not used to infer saturation.

\subsection{Two explicit algebraic source tensors}
\label{eight:sec:sources}
Each occurrence below has its own marked source copy. A negative sign
in a relation means the conjugate embedding; the coefficient $-2$ in
$r$ requires two distinct copies of $X_1$.
The following source presentations have exactly the indicated canonical
cohomological types, with nonzero entries and character order $42$:
\[
\begin{array}{c|c|c}
\text{relation}&\text{factor}&\text{Fermat triple}\\\hline
q&0&(1,1,40)\\
q,r&1&(1,13,28)\\
q&2&(1,14,27)\\
r&0&(1,20,21)\\
r&3&(1,6,35).
\end{array}
\]
Indeed $\sum_j\langle gt_j\rangle_{42}=42$ selects $\Psi_i$
for each row. Each associated rational character factor has rank-one
$M$ first cohomology and dimension six; the type-$1$ factor need not
be simple. The marked isogeny argument above identifies it with the
required $X_i$. The source embeddings for $q$ are
$(0,41),(0,25),(1,1),(2,1)$; those for $r$ are
$(0,1),(0,37),(0,13),(1,41),(1,41),(3,1)$, with the repeated
pair occupying different copies.

Concatenate the four triples after their indicated multipliers for $q$.
The resulting multiset has the following partition:
\[
 (1,41),\quad(1,41),\quad(2,13,25,25,27,34),\quad(14,28).
\]
For the six source copies of $r$ a partition is
\[
\begin{gathered}
 (1,41),\quad(1,41),\quad(6,20,21,37),\\
 (8,21,26,29),\quad(13,29),\quad(14,14,21,35).
\end{gathered}
\]
Every block $C$ is balanced:
$\sum_{a\in C}\langle ga\rangle_{42}=21|C|$ for every $g\in G$.
The only length-six block splits into the two nonzero zero-sum triples
$(2,13,27)$ and $(25,25,34)$. Thus the blocks satisfy exactly
Shioda's length-two and length-four conditions and, for length six,
his additional product-of-triples condition
\cite[Theorem~4.3 and Lemmas~4.1--4.2, pp.~70--72]{Shioda1981}.
That theorem has no odd-degree restriction; it applies at level $42$.
We use its proof supportwise for these two characters, rather than
asserting that every level-$42$ invariant support has such a partition.
The permitted block eigenspaces are algebraic, and the algebraic
induction correspondences combine them into the concatenated line.

The algebraic correspondences in those lemmas transfer the supported
Fermat-character cycles to the nonzero tensor components of the marked
curve products. For $q$ the concatenation has twelve entries, hence
Fermat dimension ten and Hodge type $(5,5)$; the corresponding curve
tensor has type $(2,2)$. For $r$ the Fermat dimension is sixteen and the curve tensor has
cohomological degree six, with types $(8,8)$ and $(3,3)$ respectively. These are the required
Tate shifts, not degree-preserving identifications. Passing through
Jacobian $H^1$, the rational character projectors, and the marked
isogenies gives the desired nonzero singleton-copy tensors on the
$X_i$. All projectors are algebraic correspondences with coefficients
in the splitting field $M$. Distinct source copies keep repeated
characters in different exterior positions. Nonzero transfer follows
from the finite-degree dominant maps in Lemma~4.2: their pullback of
a marked summand is a sum of products of the curve-character lines,
and pushforward divided by the map degree recovers the original class.
The split line and the algebraic subspace are $M$-defined; inclusion
after extension to $\C$ descends faithfully to $M$. An individual
nonrational eigenline is therefore not being called a rational cycle.
The exact multiset, all twelve balance tests, actual CM types and the
six-block splitting are independently checked by
\path{experiments/all_eight_v109/validate_witness.py}.

\subsection{Integral combinations, contractions and rational descent}
\begin{proof}[Proof of Theorem~\ref{eight:thm:all}]
Work temporarily with cycle-class coefficients extended from $\Ql$ to
$M$. Rational polynomials in multiplication-by-$n$ pullback isolate
$H^1$; actual $M$-endomorphisms then give the individual embedding-line
projectors. They distinguish all actual lines, even when their Hodge
weights agree. Addition pullback and these projectors separate any
source wedge into a specified nonzero singleton-copy tensor.
Conversely diagonal pullback assembles singleton factors into any
legitimate target exterior monomial.

We need nonzero contractions for actual conjugate embedding pairs.
Choose an ample rational divisor $\theta$ on each sixfold $X_i$.
On two singleton copies the operation
\[
 (x,y)\longmapsto\int_{X_i}x\smile y\smile\theta^5
\]
is algebraic: diagonal pullback, cup product and proper pushforward.
It is nonzero on $L_{i,a}\otimes L_{i,-a}$ by the polarization's
Riemann bilinear positivity applied to $v$ and $\overline v$.
This remains valid for induced types with repeated weights;
nondegeneracy of the whole pairing alone would not establish that
individual assertion. No inverse-Lefschetz correspondence or principal
polarization is needed.

Suppose algebraic singleton tensors represent signed relations
$v_j$ and their negatives. For an integer expression
$z=\sum_j n_jv_j$, take external products, choosing $-v_j$ when
$n_j<0$. In each actual coordinate contract all opposite pairs.
The surviving tensor has the signed multiplicities of $z$ and is a
nonzero scalar multiple of its canonical tensor. It remains algebraic
with coefficients in $M$. This proves closure under arbitrary signed
integer combinations, without a conformal decomposition or extraction
of roots of cycle classes.

The known classes and the two source tensors therefore give their
Galois translates. Here coefficient Galois action preserves the
rational cycle-class image and translates all embedding labels
simultaneously. An actual translate need not equal $J$ of its reduced
translate, but their difference lies in $\ker P$. The sixteen divisor
relations and the contraction argument correct this difference.
The integer certificate then supplies every reduced kernel relation,
and the split exact sequence supplies every actual kernel relation.

Finally take an invariant exterior monomial on any requested product.
Remove matched occurrences of opposite actual labels, accounting for
them by divisor classes. The remaining signed vector belongs to
$\ker_{\Z}W_{\rm act}$ and hence has an algebraic singleton tensor.
Restore the removed pairs and assemble the tensor into the target
copies by diagonals. No duplicate identical embedding appears within
one target exterior factor, so the result is nonzero. Auxiliary source
types absent from the target are contracted away completely.
All invariant split monomials are thus algebraic. They span the
extension to $M$ of the rational Hodge space; faithful scalar extension
implies equality of its rational subspace with the rational cycle-class
image. Complex isogeny invariance gives the stated arbitrary factors.
\end{proof}

This conclusion is limited to rank-one first cohomology over the fixed
field $M$. It neither asserts the Hodge conjecture for all CM fields
nor identifies any new geometric McMullen point, polarized isomorphism,
or principally polarized moduli normalization.

\begin{corollary}[CM factors with fields contained in the fixed field]
\label{eight:cor:subfields}
Let $A$ be a complex abelian variety. Suppose that every simple factor $B$
in a complex isogeny decomposition of $A$ admits an embedding
$E\hookrightarrow\End^0(B)$ of a CM field with $[E:\Ql]=2\dim B$,
and that $E$ admits a $\Ql$-field embedding into
$M=\Ql(\zeta_{21})$. Then every rational Hodge class on $A$ is algebraic.
The same holds for arbitrary finite mixed products of varieties satisfying
this hypothesis, including all their powers.
\end{corollary}
\begin{proof}
Write $A\sim\prod_j B_j^{m_j}$ over $\C$, and identify each chosen
$E_j$ with its image in $M$. Set $e_j=[M:E_j]$ and $X_j=B_j^{e_j}$.
Then $\dim X_j=6$. If $V_j=H^1(B_j,\Ql)$, the full-degree CM action
makes $V_j$ rank one over $E_j$, and an $E_j$-basis of $M$ identifies
\[
 H^1(X_j,\Ql)=V_j^{e_j}\simeq M\otimes_{E_j}V_j.
\]
Multiplication on the first factor gives a faithful rank-one $M$-action
preserving the Hodge decomposition. Its type is the inverse image of the
$E_j$-type under restriction of embeddings. The weight-one lattice argument
in the proof of Theorem~\ref{six:thm:products} realizes this action in
$\End^0(X_j)$; maximal-order integrality on the original lattice is not
required. Thus every $X_j$ meets Theorem~\ref{eight:thm:all}.

Choose $N_j$ with $e_jN_j\geq m_j$. Theorem~\ref{eight:thm:all} gives
rational Hodge algebraicity on
$Z=\prod_j X_j^{N_j}=\prod_j B_j^{e_jN_j}$.
Projection onto $A'=\prod_j B_j^{m_j}$ has an algebraic section inserting
origins in the extra factors. For a Hodge class $\alpha$ on $A'$, pullback
to $Z$ is algebraic; pullback along the section recovers $\alpha$ and
preserves algebraic cycle classes. Complex isogeny invariance transfers
this conclusion to $A$. A mixed product of varieties satisfying the
hypothesis has simple factors satisfying the same hypothesis, so the
argument also proves the last assertion. We use the specific mixed-product
theorem above, not a general closure of the Hodge conjecture under products.
\end{proof}
\begin{corollary}[An intrinsic splitting-field criterion]
\label{eight:cor:mt-splitting}
Let $M=\Ql(\zeta_{21})$ and let $A$ be a complex abelian variety.
The following conditions are equivalent:
\begin{enumerate}
\item $\operatorname{MT}(H_1(A,\Ql))$ is a $\Ql$-torus split over $M$;
\item every simple isogeny factor $B$ of $A$ admits a unital embedding
of a CM field $E$ into $\End^0(B)$, with $[E:\Ql]=2\dim B$, and a
$\Ql$-field embedding $E\hookrightarrow M$.
\end{enumerate}
Every rational Hodge class on an abelian variety satisfying these
conditions is algebraic. The same holds for every finite mixed product
of such varieties, including all their powers.
\end{corollary}
\begin{proof}
We use homology covariantly: $H_1(B,\Ql)$ is a polarizable Hodge structure
of weight $-1$ and type $(-1,0),(0,-1)$. The equivalence between these
rational Hodge structures and complex abelian varieties up to isogeny
identifies their endomorphisms; see \cite[Theorem~6.8 and Corollary~6.9]{MilneShimura2017}.

Suppose first that (1) holds. For a simple factor $B$, put
$V=H_1(B,\Ql)$ and let $T$ be the image of $\operatorname{MT}(A)$ on $V$.
This is a torus split over $M$ and contains the Hodge homomorphism.
The rational Hodge structure $V$ is simple: a proper nonzero Hodge
substructure would give a proper abelian subfactor up to isogeny.
Consequently $V$ is an irreducible rational representation of $T$.

Let $R\subset\End_{\Ql}(V)$ be the rational algebra generated by this
representation. Equivalently, its scalar extension contains the geometric
image of $T$; rational points need not be used. Over $M$, distinct
characters decompose $V_M$ into weight spaces, and $R_M$ is the full
algebra of scalars on those distinct spaces. Indeed, distinct characters
are linearly independent as regular functions, so their evaluations span
the corresponding diagonal algebra. Thus $R$ is finite \'etale and
commutative. Its invariant rational subspaces are exactly those of $T$.
Since $V$ is a faithful simple $R$-module, $R$ is a field $E$ and $V$ has
dimension one over $E$. Therefore $[E:\Ql]=2\dim B$, while
$E\otimes_{\Ql}M\simeq M^{[E:\Ql]}$ gives an embedding $E\hookrightarrow M$.
The algebra $E$ commutes with the Hodge homomorphism, so the covariant
Hodge equivalence embeds it into $\End^0(B)$.

It remains to see that $E$ is CM. As a subfield of the abelian Galois
extension $M/\Ql$, it is preserved by complex conjugation. If conjugation
acted trivially, $E$ would be totally real. The complex structure $J$ on
$V_{\mathbf R}$ commutes with $E$, and rank one would then give
\[
 J\in\End_{E\otimes\mathbf R}(V_{\mathbf R})
       =E\otimes\mathbf R\simeq\mathbf R^{[E:\Ql]},
\]
contradicting $J^2=-1$. Conjugation is therefore nontrivial; its fixed
field is totally real and $E$ is a totally imaginary quadratic extension
of that field. This proves (2).

Conversely, the field in (2) makes $H_1(B,\Ql)$ rank one over $E$.
The Hodge homomorphism lies in the centralizer of $E$, namely
$\operatorname{Res}_{E/\Ql}\mathbf G_m$. Minimality of the Mumford--Tate
group puts $\operatorname{MT}(B)$ inside this torus. The Mumford--Tate
group is connected: the connected Hodge torus lands in its identity
component, and minimality excludes further components. Thus it is a
subtorus. Since $M/\Ql$ is Galois and contains $E$, every embedding of
$E$ takes values in $M$, so this restriction-of-scalars torus, and each
of its subtori, splits over $M$.
The Mumford--Tate group of a finite direct sum embeds into the product
of the individual groups. It is consequently a subtorus split over $M$;
equality with that product is unnecessary. Isogenies and repetitions
preserve the conclusion, proving (1).

Corollary~\ref{eight:cor:subfields} now gives Hodge algebraicity.
The same subtorus argument proves that arbitrary finite mixed products
again satisfy (1), giving the last assertion.
\end{proof}
Here the Mumford--Tate group must already be a torus over $\Ql$.
Splitting over $M$ is an additional condition on that torus, not a claim
that a noncommutative group becomes a torus after extending scalars.

